\subsection{紧算子}\label{sec:compact-operators}\label{sec:integral-operators}

\begin{Definition}
\label{def:compact-operator}
设 $X$ 和 $Y$ 为 Banach 空间, $(\mathcal{D}(A),A)$ 为从 $X$ 到 $Y$ 的线性算子. 若算子 $A$ 将有界集映为预紧集, 即对每个有界集 $\Omega\subset\mathcal{D}(A)$, 都有 $\overline{A(\Omega)}\subset Y$ 为紧集, 则称它为\textbf{紧算子}.
\end{Definition}

用序列而非集合来刻画紧算子往往更方便.

\begin{Theorem}
\label{thm:compact-operator-sequential-characterization}
从 $X$ 到 $Y$ 的算子 $(\mathcal{D}(A),A)$ 是紧的, 当且仅当它是序列紧的; 即当且仅当对 $\mathcal{D}(A)$ 中任意有界序列 $x_n$, $A(x_n)$ 都有收敛子序列.
\end{Theorem}

\begingroup
\postdisplaypenalty=10000
\renewcommand{\endProof}{\par\nobreak\hfill$\blacksquare$\par\vspace{0.1cm}}
\begin{Proof}
该定理的证明直接来自预紧集可以用序列刻画这一拓扑结果; 即赋范线性空间中的集合 $S$ 预紧, 当且仅当包含在 $S$ 中的每个序列都有收敛子序列.
\end{Proof}
\endgroup

正如下文将看到的, 紧算子的最基本例子是积分算子. 然而, 要更充分地研究它们, 还需要建立一些工具. 与此同时, Section~\ref{subsec:sobolev-compactness-properties} 中对紧嵌入的研究已经提供了一些非常重要的紧算子例子. 为解释这些例子, 需要以下引理.

\begin{Lemma}
\label{lem:compact-embedding-identity-operator}
设 $X$ 和 $Y$ 为 Banach 空间. 那么 $X$ 紧嵌入 $Y$, 当且仅当从 $X$ 到 $Y$ 的恒等映射有良好定义且为紧映射.
\end{Lemma}

由 Definition~\ref{def:compact-imbedding} 以及 Example~\ref{ex:identity-operator-and-inclusion} 中恒等映射的定义, 立即得到证明.

\begin{Example}
\label{ex:sobolev-continuous-target-compact-operator}
由 Theorem~\ref{thm:sobolev-continuous-target-compact-embedding} 可知, 若 $k>m/2$, $\Omega\subset\mathbb{R}^m$ 有界且具有 $k$ 阶延拓性质, 则从 $H^k(\Omega)$ 到 $C_b(\overline\Omega)$ 的恒等映射为紧映射. 因而由 Theorem~\ref{thm:compact-operator-sequential-characterization}, 每个在 $H^k(\Omega)$ 范数下有界的函数序列 $u_n$ 都有一致收敛的子序列.
\end{Example}

\begin{Example}
\label{ex:successive-sobolev-orders-compact-operator}
由 Theorem~\ref{thm:compact-successive-sobolev-orders} 可知, 若 $k$ 为非负整数, $\Omega\subset\mathbb{R}^m$ 有界且具有 $k+1$ 阶延拓性质, 则从 $H^{k+1}(\Omega)$ 到 $H^k(\Omega)$ 的恒等映射为紧映射. 再次利用 Theorem~\ref{thm:compact-operator-sequential-characterization}, 可知每个在 $H^{k+1}(\Omega)$ 范数下有界的函数序列 $u_n$, 都有在 $H^k(\Omega)$ 范数下强收敛的子序列.
\end{Example}

现在得到以下初等结果.

\begin{Lemma}
\label{lem:compact-operator-boundedness}
每个紧算子都有界.
\end{Lemma}

\begingroup
\postdisplaypenalty=10000
\renewcommand{\endProof}{\par\nobreak\hfill$\blacksquare$\par\vspace{0.1cm}}
\begin{Proof}
反设不然, 则存在序列 $x_n\in\mathcal{D}(A)$, 使得 $\|x_n\|=1$ 且 $\|A(x_n)\|\to\infty$. 事实上, 通过删除序列中的多余元素并重新编号, 可保证 $\|A(x_{n+1})\|>\|A(x_n)\|+1$. 因而 $A(x_n)$ 的任何子序列都不可能收敛, 因为没有子序列可能是 Cauchy 序列.
\end{Proof}
\endgroup

回忆 Theorem~\ref{thm:bounded-operator-norm-preserving-extension} 表明, 每个有界算子都可扩张到整个 $X$, 而不改变其范数. 留给读者证明, 用 Theorem~\ref{thm:bounded-operator-norm-preserving-extension} 的证明中所述的方法扩张紧算子时, 扩张后的算子仍为紧算子 (Problem~\ref{ex:compact-operator-norm-preserving-extension}). 因而通常假设紧算子属于 $\mathcal{L}(X,Y)$.

注意, Lemma~\ref{lem:compact-operator-boundedness} 和 Lemma~\ref{lem:linear-mapping-bounded-continuity-equivalence} 表明每个紧算子都连续. 然而, 这一结果的逆命题不成立. 特别地, 有以下结果.

\begin{Lemma}
\label{lem:infinite-dimensional-identity-noncompact}
若 $X$ 为任意无穷维 Banach 空间, 则恒等算子不是紧算子.
\end{Lemma}

\begingroup
\postdisplaypenalty=10000
\renewcommand{\endProof}{\par\nobreak\hfill$\blacksquare$\par\vspace{0.1cm}}
\begin{Proof}
证明直接来自无穷维空间中单位球不紧这一事实. 这里只证明无穷维 Hilbert 空间的情形, 一般结果留给读者 (Problem~\ref{ex:infinite-dimensional-identity-noncompact}). 回忆 Corollary~\ref{cor:hilbert-basis-existence} 表明, 无穷维 Hilbert 空间中存在无穷标准正交集 $\{x_i\}_{i=1}^\infty$. 该集合包含于闭单位球, 而对基中两个不同元素 $x_i$ 和 $x_j$, 有 $\|x_i-x_j\|^2=2$. 因而 $x_i$ 的任何子序列都不可能收敛, 因为没有子序列可能是 Cauchy 序列.
\end{Proof}
\endgroup
